\titleformat
{\chapter} % command
[display] % shape
{\bfseries\Large\itshape} % format
{Capítulo \ \thechapter} % label
{0.5ex} % sep
{
    \rule{\textwidth}{1pt}
    \vspace{1ex}
    \centering
} % before-code
[
\vspace{-0.5ex}%
\rule{\textwidth}{0.3pt}
] % after-code

\chapter{Planteamiento del Problema}

\section{Introducción al Problema}

El uso acelerado de dispositivos electrónicos ha incrementado drásticamente la generación de residuos de aparatos eléctricos y electrónicos (RAEE). Estos residuos contienen sustancias tóxicas que afectan el medio ambiente y la salud humana, agravados por la obsolescencia programada y bajos niveles de reciclaje a nivel global \cite{jesiddiazbastosImpactoResiduosAparatos2025}. En América Latina, la inadecuada gestión de esta producción masiva genera una acumulación crítica que exige la implementación de mejores prácticas de manejo \cite{zambranoyepezBuenasPracticasManejo2022}.

\section{Análisis del Problema}

Los RAEE contienen materiales valiosos (oro, cobre, aluminio), representando una fuente potencial de recuperación económica. Sin embargo, su compleja caracterización y la ineficiencia operativa limitan su aprovechamiento \cite{venegasmarcelModeloProcedimentalPara2020}. En países en desarrollo, la falta de infraestructura y de procesos avanzados provoca la exportación de estos residuos, transfiriendo su valor económico y reduciendo la viabilidad de soluciones locales \cite{ottoniSustainabilityAssessmentBest2023}. Esta gestión ineficiente dificulta la implementación de modelos de economía circular, incrementando costos logísticos y desperdiciando materiales recuperables \cite{souzaCHALLENGESTRANSITIONCIRCULAR2025}.\\

En el contexto colombiano, la gestión de los RAEE ha sido fortalecida mediante lineamientos técnicos y marcos normativos orientados a su recolección, tratamiento y disposición final. Sin embargo, a pesar de estos avances, persisten limitaciones relacionadas con la baja adopción de tecnologías en entornos operativos, la dependencia de procesos manuales en la clasificación y la falta de articulación entre las etapas de identificación y valorización de los residuos, lo que limita el aprovechamiento eficiente de los materiales recuperables y evidencia brechas en la implementación de soluciones tecnológicas integradas \cite{guiaRAEEColombia2023}.\\

Con el fin de sintetizar las causas y efectos asociados a la problemática identificada, en la Figura \ref{fig:arbol_problema_final} se presenta el árbol del problema.\\

\begin{figure}[H]
\centering
    \resizebox{\textwidth}{!}{
        \begin{tikzpicture}[
            node distance = 0.5cm and 0.5cm,
            base/.style = {rectangle, draw, line width=1pt, text width=4cm, minimum height=1.6cm, align=center, font=\sffamily\scriptsize},
            efecto/.style = {base, draw=green!60!black, fill=green!5},
            problema/.style = {base, draw=orange!70!black, fill=orange!10, text width=5.5cm, minimum height=2.8cm},
            causa/.style = {base, draw=brown!70!black, fill=brown!15},
            labelStyle/.style = {font=\sffamily\tiny\bfseries, rotate=90}
        ]

        % PROBLEMA
        \node[problema] (P) {Dificultad para clasificar y valorar económicamente los RAEE en tiempo real debido a la dependencia de procesos manuales y la limitada integración de tecnologías basadas en inteligencia artificial.};

        % EFECTOS
        \node[efecto, above=1.0cm of P] (E2b) {Riesgos para la salud en la manipulación de residuos peligrosos.};
        \node[efecto, left=of E2b] (E1b) {Impacto ambiental por disposición inadecuada de RAEE.};
        \node[efecto, right=of E2b] (E3b) {Baja competitividad en el mercado de materiales reciclables.};

        \node[efecto, above=of E1b] (E1a) {Pérdida de valor económico en materiales reciclables.};
        \node[efecto, above=of E2b] (E2a) {Incremento de costos operativos en los procesos de clasificación.};
        \node[efecto, above=of E3b] (E3a) {Baja eficiencia en el aprovechamiento de los residuos.};

        % CAUSAS
        \node[causa, below=1.0cm of P] (C2a) {Baja adopción de tecnologías de inteligencia artificial en entornos operativos.};
        \node[causa, left=of C2a] (C1a) {Dependencia de inspección visual manual en el proceso de clasificación.};
        \node[causa, right=of C2a] (C3a) {Falta de integración entre clasificación visual y valoración económica.};

        \node[causa, below=of C1a] (C1b) {Infraestructura tecnológica limitada y de alto costo para el procesamiento de RAEE.};
        \node[causa, below=of C2a] (C2b) {Alta complejidad en la caracterización de materiales electrónicos.};
        \node[causa, below=of C3a] (C3b) {Desconexión entre clasificación en planta y comercialización de materiales.};

        % CONEXIONES
        \begin{scope}[gray, thin]

            \draw (P.north) -- ++(0,0.5) -| (E1b.south);
            \draw (P.north) -- (E2b.south);
            \draw (P.north) -- ++(0,0.5) -| (E3b.south);

            \draw (P.south) -- ++(0,-0.5) -| (C1a.north);
            \draw (P.south) -- (C2a.north);
            \draw (P.south) -- ++(0,-0.5) -| (C3a.north);

            \draw (E1b.north) -- (E1a.south); \draw (E2b.north) -- (E2a.south); \draw (E3b.north) -- (E3a.south);
            \draw (C1a.south) -- (C1b.north); \draw (C2a.south) -- (C2b.north); \draw (C3a.south) -- (C3b.north);
        \end{scope}

        % LLAVES
        \draw [decorate, decoration={brace, amplitude=10pt, mirror}, green!40!black, line width=1pt]
            ([xshift=-10pt]E1a.north west) -- ([xshift=-10pt]E1b.south west)
            node [midway, left=15pt, labelStyle, anchor=center] {EFECTOS};

        \draw [decorate, decoration={brace, amplitude=10pt, mirror}, orange!60!brown, line width=1pt]
            ([xshift=-10pt]P.north west) -- ([xshift=-10pt]P.south west)
            node [midway, left=15pt, labelStyle, anchor=center] {PROBLEMA};

        \draw [decorate, decoration={brace, amplitude=10pt, mirror}, brown!60!black, line width=1pt]
            ([xshift=-10pt]C1a.north west) -- ([xshift=-10pt]C1b.south west)
            node [midway, left=15pt, labelStyle, anchor=center] {CAUSAS};

        \end{tikzpicture}
    }
\caption{Árbol del problema para la gestión de RAEE}
\label{fig:arbol_problema_final}
\end{figure}

\vspace{1em}

Ante estas limitaciones, los enfoques basados en inteligencia artificial evidencian un alto potencial para automatizar la clasificación de residuos. Las técnicas de visión computacional y aprendizaje profundo, específicamente las redes neuronales convolucionales (CNN) y arquitecturas como YOLO, permiten identificar defectos y clasificar dispositivos con alta precisión en tiempo real \cite{menagaIntelligentDeviceCondition2026}. Mejoras como YOLO-MDD abordan la variabilidad de escalas y condiciones de captura, superando a modelos tradicionales en métricas de robustez \cite{guYOLOMDDAddressingInconsistent2025}.\\

Adicionalmente, se ha explorado la integración de visión computacional con esquemas de valoración económica para estimar el valor potencial del residuo según su estado \cite{joshiSustainableEWasteManagement2024}, así como el uso de enfoques multimodales para reducir la ambigüedad en clasificaciones complejas \cite{obayemiAutomatedWEEESortation2025a}. Por su parte, arquitecturas recientes como los Vision Transformers alcanzan precisiones superiores al 99\% \cite{chauhanViT4WasteVisionTransformerBased2025}. Sin embargo, su alto costo computacional, consumo de memoria y elevada latencia limitan drásticamente su aplicabilidad en escenarios operativos, haciendo imperativo el uso de redes eficientes orientadas al procesamiento continuo.\\

En contraste, la literatura reciente evidencia que los modelos ligeros pueden ejecutarse en plataformas de cómputo de bajo costo. \cite{castrobelloConvolutionalNeuralNetwork2025} compararon cuatro detectores ligeros (YOLOv4-tiny, YOLOv7-tiny, YOLOv8-nano y YOLOv9-tiny) sobre una Raspberry Pi 4 y reportaron un mAP50 de 0{,}976 para YOLOv8-nano, con tiempos de inferencia cercanos a 1{,}5\,s por imagen; \cite{daoDesignImplementationRealTime2026} desplegaron YOLOv8n en la CPU de una Raspberry Pi 5 con una latencia de 150\,ms ($\sim$6{,}7\,FPS), y \cite{mutarAIBasedModelHome2026} alcanzaron un \textit{accuracy} de 95{,}2\% con 45\,ms por imagen y un consumo promedio de 5{,}2\,W al incorporar el acelerador neuronal del Raspberry Pi 5 AI Kit. Estos resultados indican que la barrera para la adopción en centros de acopio no es la disponibilidad de algoritmos, sino su adaptación a infraestructuras de cómputo asequibles, condición todavía poco explorada para la clasificación de componentes de RAEE.\\

A pesar de estos avances algorítmicos, su implementación real presenta barreras críticas. La recuperación de materiales sigue siendo parcial y fuertemente dependiente de procesos manuales, lo que genera grandes pérdidas de valor económico \cite{oliveiranetoCriticalBarriersMaterial2025}. La complejidad del manejo de RAEE representa un problema de decisiones multicriterio, gravemente afectado por ineficiencias logísticas y altos costos de operación \cite{seikhIdentifyingSustainableStrategies2024, oliveiranetoArtificialIntelligencebasedOptimization2025}.\\

Las soluciones actuales suelen enfocarse en silos aislados, como la optimización de rutas o la recuperación final, sin integrar de manera efectiva la clasificación visual con modelos de valoración financiera \cite{simaeiAIbasedDecisionSupport2024}. Esta falta de información estructurada y la alta incertidumbre en la condición del residuo dificultan la toma de decisiones en etapas clave como la reutilización, reparación o reciclaje \cite{anandhMultiobjectiveMetaheuristicsOptimize2024}.\\

Por lo tanto, se propone el desarrollo de un sistema de apoyo a la toma de decisiones basado en técnicas de visión computacional y un esquema de valoración económica, ejecutable en hardware de cómputo de bajo costo. Este sistema permitirá identificar y clasificar residuos electrónicos a partir de imágenes, estimando simultáneamente su potencial de rentabilidad para reducir tiempos operativos y optimizar el proceso de reciclaje, sin requerir infraestructura computacional especializada en los centros de acopio.

\section{Formulación del Problema}
¿En qué medida la implementación de un sistema de apoyo a la toma de decisiones, que integre modelos de visión computacional con un esquema de valoración económica y que se ejecute en hardware de cómputo de bajo costo, permite mejorar la precisión, reducir los tiempos de clasificación y optimizar la viabilidad del proceso en la gestión de RAEE frente a métodos tradicionales?
